\chapter*{A Note on Structure}
\addcontentsline{toc}{chapter}{A Note on Structure}

This book intentionally presents several frameworks that are later revised, corrected, or abandoned. The first model of elimination (Chapter~3) turns out to be unrepairable and is replaced in Chapter~4. The first account of how a child's vocabulary develops (Chapter~7) turns out to be wrong about what a ``subsystem'' is, and is repaired mid-chapter. A clean curriculum hypothesis in Chapter~11, and a clean independence hypothesis in Chapter~12, both meet the same fate. This is not because the manuscript was undecided about its own claims while being written, and a reader encountering the pattern cold could be forgiven for suspecting otherwise.

The central claim of this book is that concepts become trustworthy not by being declared correct, but through repeated cycles of proposal, criticism, elimination, and repair --- and that the difference between a hard core and a soft periphery only becomes visible once something has been tested against reality and found wanting. A book that only ever showed its readers the final, successful version of each idea would be describing that process from the outside. This one tries, instead, to put the reader inside it: each provisional model is built in earnest, allowed to fail on its own terms, and then repaired in a way the reader can check against the failure that prompted it.

To make this easier to track while reading, provisional models that are later revised are marked in a shaded box like this one is not, labeled \textsc{provisional model}, and the repair that replaces each one is marked in a second, differently shaded box labeled \textsc{repair}. A box of either kind signals exactly what it is: a stage in an argument, not the book's last word on the subject. The book's actual closing position on any given question is always the content of the final repair box in the relevant chapter, or, failing that, the surrounding prose once the boxes stop appearing.

One further piece of apparatus is marked differently still. Section~17.4 builds a small amount of physics-flavored formal machinery --- a Lagrangian, an Euler--Lagrange equation, a Hamiltonian --- to describe, speculatively, how the book's elimination dynamics might behave as a continuous process. This material is explicitly conjecture-level, clearly subordinate to the operational definitions that precede it in the same section, and is marked as an \textsc{optional formal extension}: a reader uninterested in the dynamical apparatus can skip the boxed material entirely without losing anything the rest of the book depends on.

Why this structure was chosen, rather than the more conventional alternative of presenting only the finished framework, is explained directly in the book's final chapter, once the reader has the full pattern in view to judge for themselves whether it worked.
